\documentclass[a4paper,12pt]{article}
\usepackage[utf8]{inputenc}
\usepackage[heading=true, fontset=mac]{ctex}
\usepackage{geometry}
\usepackage{graphicx}
\usepackage{float}
\usepackage{fancyhdr}
\usepackage{hyperref}
\usepackage{setspace}
\usepackage{titlesec}
\usepackage{xcolor}

% 页面设置
\geometry{left=2.5cm, right=2.5cm, top=2.5cm, bottom=2.5cm}
\setlength{\parindent}{2em}
\onehalfspacing

% 页眉页脚配置
\pagestyle{fancy}
\fancyhf{}
\fancyhead[C]{[项目名称] 专利交底书 (Internal Ref: [编号])}
\fancyfoot[C]{\thepage}

% 标题信息
\title{\textbf{[发明名称：一种基于...的方法/系统]}}
\author{发明人：AI4S 团队}
\date{\today}

\begin{document}

\maketitle

\section{著录项目}
\begin{itemize}
    \item \textbf{发明名称：} [全称]
    \item \textbf{技术领域：} [具体领域1]、[具体领域2]
    \item \textbf{关键词：} [关键词1]，[关键词2]，[关键词3]
\end{itemize}

\section{背景技术与技术问题}

\subsection{现有技术}
\textit{（简述当前主流的解决方案，例如 SFT, RLHF, 传统数值求解器等）}
\begin{enumerate}
    \item \textbf{[方案一]：} ...
    \item \textbf{[方案二]：} ...
\end{enumerate}

\subsection{技术缺陷}
\textit{（客观描述现有技术无法解决的具体痛点，避免主观贬低）}
现有技术在处理 [特定场景] 时，存在以下缺陷：
\begin{itemize}
    \item \textbf{缺陷一（[命名]）：} ...
    \item \textbf{缺陷二（[命名]）：} ...
\end{itemize}

\begin{figure}[H]
    \centering
    \fbox{\parbox{0.8\textwidth}{\centering \vspace{1.5cm} [示意图 1：现有技术的缺陷/问题场景] \vspace{1.5cm}}}
    \caption{现有技术的问题示意图}
    \label{fig:problem}
\end{figure}

\subsection{技术问题}
本发明要解决的核心技术问题是：\textbf{如何在 [约束条件] 下，实现 [核心目标]（例如：逻辑一致性、物理守恒性）。}

\section{技术方案}

\subsection{总体思路}
本发明提出了一种 \textbf{[架构名称]}。其核心思路在于 [一句话核心机制]（例如：引入可微一致性层，将逻辑约束转化为训练损失）。

\subsection{与现有技术的区别}
\begin{itemize}
    \item \textbf{对比 [现有技术A]：} [现有技术A] 侧重于 ...，而本发明侧重于 ...（本质区别）。
    \item \textbf{机制差异：} 本发明通过 [具体手段] 避免了 [常见问题]（例如：无需重训练、解耦设计）。
\end{itemize}

\subsection{系统架构图}
\begin{figure}[H]
    \centering
    \fbox{\parbox{0.9\textwidth}{\centering \vspace{2cm} [示意图 2：本发明的总体系统架构图] \vspace{2cm}}}
    \caption{系统架构图}
    \label{fig:architecture}
\end{figure}

\subsection{核心模块详解}

\subsubsection{模块一：[名称] ($\mathcal{M}_1$)}
\begin{itemize}
    \item \textbf{功能：} ...
    \item \textbf{输入/输出：} 输入 $X$，输出 $Y$。
    \item \textbf{处理逻辑：} ...
\end{itemize}

\begin{figure}[H]
    \centering
    \fbox{\parbox{0.8\textwidth}{\centering \vspace{1.5cm} [示意图 3：核心模块/算法流程图] \vspace{1.5cm}}}
    \caption{核心模块工作原理图}
    \label{fig:module1}
\end{figure}

\subsubsection{模块二：[名称] ($\mathcal{M}_2$)}
\begin{itemize}
    \item \textbf{功能：} ...
    \item \textbf{核心算法：} ...
\end{itemize}

\section{有益效果}

\begin{enumerate}
    \item \textbf{[效果一]：}
    通过 [技术手段]，实现了 [量化指标提升]（例如：F1 分数提升至 XX\%）。
    
    \item \textbf{[效果二]：}
    解决了 [具体缺陷]，使得 [应用场景] 成为可能。
\end{enumerate}

\begin{figure}[H]
    \centering
    \fbox{\parbox{0.6\textwidth}{\centering \vspace{1.5cm} [示意图 4：性能对比/效果展示图] \vspace{1.5cm}}}
    \caption{本发明与现有技术的性能对比}
    \label{fig:performance}
\end{figure}

\section{具体实施方式}

\subsection{实施例 1：[场景名称]}
\begin{itemize}
    \item 步骤 1：...
    \item 步骤 2：...
\end{itemize}

\section{权利要求书草案}

\begin{enumerate}
    \item 一种[发明名称]方法，其特征在于，包括：
    \begin{itemize}
        \item 步骤 S1：...
        \item 步骤 S2：...
        \item 步骤 S3：...
    \end{itemize}
    \item 如权利要求 1 所述的方法，其特征在于，所述步骤 S2 包括...
    \item 一种系统，包括处理器，配置为执行权利要求 1 所述的方法。
\end{enumerate}

\end{document}
